% !TEX root = chaosproject.tex
\include{preamble.tex}
\def\title{Lecture Notes}
\def\coursenumber{}
\def\class{Matrix Concentration and Applications}
\def\semester{}
\def\name{Rohit Agarwal}
\def\title{Notes}
\lhead{}
\chead{\title}
\rhead{\name}
\lfoot{}
\rfoot{\thepage}

\begin{document}


\newpage
\section{Matrix Chaos Inequalities}
For the purposes of this note, $\norm{\cdot}$ when applied to a matrix will always denote spectral norm.
We will present the main results of \cite{bandeira2025matrixchaosinequalitieschaos}.
In this work, they were investigating generalizations of the Matrix Khintchine, where they study a different
matrix model, called matrix chaos. In this model, the component
matrices are multiplied by a polynomial of random variables. Here, we will study
a restricted subset of these.
\begin{definition}
    Let $m$ and $q$ be integers.
    Let $\{h_{i}^{(j)}\}_{i \in [m], j \in [q]}$ be iid copies of a random variable $h$ with zero mean,
    and let $A_{i_1, \dots, i_q}$ be a set of $m^q$ deterministic $d_1 \times d_2$ matrix coefficients.
    Then, a matrix of the form
    \[ X = \sum_{i_1, \dots, i_q \in [m]} h_{i_1}^{(1)} \dots h_{i_q}^{(q)} A_{i_1, \dots, i_q}  \]
    is called a \textbf{matrix chaos} of order $q$.
\end{definition}
It's known that $\E{\norm{X}}$ does not change much (changes by $q$-dependent factors) when we replace this ``decoupled'' version with
1) the dependence that $h_{i}^{(1)} = h_{i}^{(2)} = \dots = h_{i}^{(q)}$ and 2) forbidding $i_k$ from repeating (enforcing ``square-freeness'').
As we will see, this form is what pops up in the applications and the inequalities. 

To give the bounds, we first need to define
the notion of a flattening. Fix a model $X$ as above.
Then, identify the matrix coefficients with a tensor $\mathcal{A}$ of order $q + 2$ by
\[ \mathcal{A}_{i_1, \dots, i_q, i_{q + 1}, i_{q + 2}} = A_{i_1, \dots, i_q}[i_{q + 1}, i_{q + 2}] \]
The first $q$ coordinates we will call chaos coordinates and the last two coordinates we will call matrix coordinates.
\begin{definition}
    Given a partition of $[q + 2]$ as $R \cup C$, we define the flattening of $\mathcal{A}$ as the matrix
    \[ \mathcal{A}_{[R \mid C]} = \sum_{\substack{i_1, \dots, i_q \in [m]}\\
    i_{q + 1} \in [d_1], i_{q + 2} \in [d_2]} \qty(\bigotimes_{ t\in R} e_{i_t}) \otimes \qty(\bigotimes_{t \in C} e_{i_t}^T) \mathcal{A}_{i_1, \dots, i_{q + 2}}\]
\end{definition}
Here, this is the matrix whose rows are indexed by the coordinates in the row set $R$, whose columns
are indexed by the coordinates in the column set $C$, and whose entries are the corresponding entries of the tensor.
\begin{definition}
    A flattening is said to be a $\sigma$-flattening if $(q + 1) \in R$ and $(q + 2) \in C$,
    i.e. the matrix coordinates are kept as row and column coordinates. Then,
    \[ \sigma(\mathcal{A}) := \max_{[R\mid C], \, \sigma\text{-flattening}} \norm{\mathcal{A}_{[R \mid C]}} \]
    A flattening is said to be a $v$-flattening if $(q + 1), (q + 2) \in C$ i.e. the matrix coordinates are both assigned to the columns and $R$ is nonempty.
    Then,
    \[ v(\mathcal{A}) := \max_{[R\mid C], \, v\text{-flattening}} \norm{\mathcal{A}_{[R \mid C]}}  \]
\end{definition}
We will show how to reformulate the usual NCK in terms of flattenings (i.e. $q = 1$ case),
which is the innovation of the chaos inequalities. Then, the full chaos inequalities will follow by an induction
on $q$, which we will skip.

Let $Y = \sum_{i \in [m]} h_i A_i$. Recall that the NCK Inequality \cite{Pisier_2003} states that 
\[  \mathbb{E} \norm{Y} \lesssim \norm{h}_{\psi_2} \log(d)^{1/2} \max\{\sigma_R(Y), \sigma_C(Y)\} \]
where $\norm{h}_{\psi_2}$ is the subgaussian norm of $h$ and
\[ \sigma_R(Y) := \norm{\sum_{i \in [m]} A_i^T A_i}^{1/2}, \;\;\;\;\; \sigma_C(Y) := \norm{\sum_{i \in [m]} A_i A_i^T}^{1/2} \]
Then we can compute the following flattenings
\[ \mathcal{A}_{[1, 2 \mid 3]} = \sum_{i \in [m]} e_i \otimes A_i = \begin{pmatrix}
    A_1 \\ \vdots \\ A_m
\end{pmatrix}, \;\;\;\;\; \mathcal{A}_{[2 \mid 1, 3]} = \sum_{i \in [m]} e_i^T \otimes A_i = \begin{pmatrix} A_1 & \dots & A_m \end{pmatrix} \]
which means
\[ \mathcal{A}_{[1, 2 \mid 3]}^T \mathcal{A}_{[1, 2 \mid 3]} = \sum_{i \in [m]} A_i^T A_i, \; \;\;\;\;\mathcal{A}_{[2 \mid 1, 3]} \mathcal{A}_{[2 \mid 1, 3]}^T = \sum_{i \in [m]} A_i A_i^T  \]
so therefore $\sigma_R(Y) = \norm{\mathcal{A}_{[1, 2 \mid 3]}}$ and $\sigma_C(Y) = \norm{\mathcal{A}_{[2 \mid 1,3]}}$,
which enumerates all the $\sigma$-flattenings of $\mathcal{A}$. Thus, we can reformulate as $\mathbb{E} \norm{Y} \lesssim \norm{h}_{\psi_2} \log(d)^{1/2} \sigma(\mathcal{A})$
for the $q = 1$ case.

We can now state what induction will recover:
\begin{theorem}[Iterated NCK]
    Let $Y$ be a matrix chaos. Then,
    \[ \norm{h}_{L_1}^q \sigma(\mathcal{A}) \lesssim_q \mathbb{E} \norm{Y} \lesssim \norm{h}_{\psi_2}^q \log(d + m)^{q/2} \]
\end{theorem}
One can also write the $v(Y) = \norm{Cov(Y)}^{1/2}$ parameter in the strong NCK of \cite{Bandeira_2023}
by noticing
\[ \mathcal{A}_{[1, \mid 2, 3]} = \sum_{i \in [m]} e_i \otimes vec(A_i)^T = \begin{pmatrix}
    vec(A_i)^T \\ \vdots \\ vec(A_m)^T
\end{pmatrix} \]
where vec denotes the operation that arranges the matrix into a long vector.
Then, one can easily notice that $\mathcal{A}^T_{[1 \mid 2, 3]} \mathcal{A}_{[1 \mid 2,3]} = Cov(Y)$,
so $v(Y) = \norm{\mathcal{A}_{[1 \mid 2, 3]}}$. This is gives rise to the following inequality.
\begin{theorem}[Iterated Strong NCK]
    Let $Y$ be a matrix chaos. Then,
    \[ \mathbb{E} \norm{Y} \lesssim \norm{h}_{\psi_2}^q(\sigma(\mathcal{A}) + \log(d + m)^{\frac{q + 2}{2}} v(\mathcal{A})) \]
\end{theorem}

\subsection{Chaos of Combinatorial Type}
Although we capture a large class of models, computing the flattening quantities can be a bit daunting.
However, the authors note that for many natural models, there is additional structure that one
can exploit to use the chaos inequalities. Here, we use them for Khatri-Rao matrices, which
have been used for differential privacy \cite{krmatrices}, among other applications.

\begin{definition}
    Let $q, n, d$ be positive integers, let $h$ be a scalar random variable with zero mean and unit variance, and let $W^{(1)}, \dots, W^{(q)}$
    be $d \times n$ random matrices whose entries are iid copies of $h$. The Khatri-Rao matrix $Y$ is the $d^q \times n$
    random matrix obtained by taking the column-wise Kronecker product of $W^{(1)}, \dots, W^{(q)}$, that is where 
    for $j_1, \dots, j_q \in [d]$ and $k \in [n]$ the entries are:
    \[ Y[(j_1, \dots, j_q), k] = \prod_{t = 1}^{q} W^{(t)}[j_t, k] \]
    Equivalently,
    \[ Y = \sum_{j_1, \dots, j_q \in [d], k\in[n]} \qty(\prod_{t = 1}^{q} W^{(t)}[j_t, k]) e_{j_1} \otimes \dots \otimes e_{j_q} \otimes e_k^T \]
\end{definition}
Here, we see that Khatri-Rao matrices have the property 
that every matrix coefficient $A_{(j_1, k), \dots, (j_q, k)}$ is a tensor product of coordinate basis vectors and their transposes\footnote{Note that to capture this under the original chaos model,
we also need a bunch of $0$ matrices e.g. for the cases where $A_{(j_1, k_1), \dots, (j_q, k_q)}$ for some $k_i \neq k_j$.
We will ignore this to keep the notation easier.}.
This is called a \textbf{matrix chaos of combinatorial type}. 
This structure enables us to compute the norms of the flattenings much more easily.

Let $Y$ be a Khatri-Rao matrix with $q = 2$. 
Thus, defining $h_{(j_t, k)}^{(t)} = W^{(t)}[j_t, k]$ we can write
\[ Y = \sum_{j_1, j_2 \in [d], k \in [n]} h_{(j_1, k)}^{(1)} h_{(j_2, k)}^{(2)} e_{(j_1, j_2)} \otimes e_k^{T}\]
Then
\begin{enumerate}
    \item The chaos coordinates are $(j_1, k), (j_2, k)$
    \item The matrix coordinates are $(j_1, j_2)$ and $k$
\end{enumerate}
Let's say we wanted to compute the flattening
\begin{align*}
    \mathcal{A}_{[1, 2, 3 \mid 4]} &= \sum_{j_1, j_2 \in [d], k \in [n]} e_{(j_1, k)} \otimes e_{(j_2, k)} \otimes e_{(j_1, j_2)} \otimes e_k^T \\
    &= \sum_{j_1, j_2 \in [d], k \in [n]} e_{j_1} \otimes e_k \otimes e_{j_2} \otimes e_k \otimes e_{j_1} \otimes e_{j_2} \otimes e_k^T
\end{align*}
where we expanded via tensor products. Permuting the order of the tensor products doesn't preserve the matrix itself,
but only reorders rows and columns, so it doesn't change the singular values.
Therefore,
\begin{align*}
    \mathcal{A}_{[1, 2, 3, \mid 4]} &\simeq \qty(\sum_{j_i \in [d]} e_{j_1} \otimes e_{j_1}) \otimes \qty(\sum_{j_2 \in [d]} e_{j_2} \otimes e_{j_2}) \otimes \qty(\sum_{k \in [n]} e_k \otimes e_k \otimes e_k^T) \\
    \norm{\mathcal{A}_{[1, 2, 3, \mid 4]}} &= \norm{\sum_{j_i \in [d]} e_{j_1} \otimes e_{j_1}}_2 \cdot \norm{\sum_{j_2 \in [d]} e_{j_2} \otimes e_{j_2}}_2 \cdot \norm{\sum_{k \in [n]} e_k \otimes e_k \otimes e_k^T} \\
    &= \sqrt{d} \cdot \sqrt{d} \cdot 1 = d
\end{align*}
where the first two are the sums of $d$ orthonormal terms, and the last is an identity matrix.
If one were to compute all the possible ways of distributing the chaos coordinates $i_1, \dots, i_q$
and matrix coordinates $i_{q + 1}, i_{q + 2}$

In general,
for any matrix chaos of combinatorial type where the summation index $u$ has support $S_u$
(in our case we have $3$ summation indices, $j_1, j_2, k$, with supports $d, d, n$),
\[ \mathcal{A}_{[R \mid C]} = \bigotimes_{u = 1}^{p} \sum_{s_u \in [S_u]} \qty(e_{S_u}^{\otimes \mu_u} \otimes (e_{S_u}^T)^{\otimes \nu_u}) \]
Now, noting that by similar reasoning:
\[ \norm{\sum_{s_u \in [S_u]}\qty(e_{S_u}^{\otimes \mu_u} \otimes (e_{S_u}^T)^{\otimes \nu_u})} = \begin{cases}
    1 & \text{if } \mu_u > 0 \text{ and } \nu_u > 0 \\ 
    \sqrt{S_u} & \text{if } \mu_u > 0 \text{ xor } \nu_u > 0 \\
    S_u & \text{if } \mu_u = 0 \text{ and } \nu_u = 0
\end{cases}  \]
Thus, we can compute all the flattenings for Khatri-Rao matrices: 
\begin{table}[hbt!]
    \centering
    \begin{tabular}{| c | c | c | c | c |}
        \hline
        \multicolumn{4}{|c|}{coordinates} &  \\
        $j_1 k$ & $j_2 k$ & $j_1 j_2$ & $k$ & $\norm{\mathcal{A}_{[R|C]}}^2$ \\ \hline
        $R$ & $R$ & $R$ & $C$ & $d^2$ \\ \hline 
        $R$ & $C$ & $R$ & $C$ & $d$ \\ \hline
        $C$ & $R$ & $R$ & $C$ & $d$ \\ \hline
        $C$ & $C$ & $R$ & $C$ & $n$ \\ \hline
    \end{tabular}
\end{table}
\\
This shows that for $q = 2$ Khatri-Rao matrices, $\sigma(\mathcal{A}) = \max\{d, n^{1/2}\}$.
Furthermore, one can repeat such a procedure on $v$ to get that $v(\mathcal{A}) = d^{1/2} \ll \sigma(\mathcal{A})$, so Iterated strong NCK gives
\[ \mathbb{E} \norm{Y} \approx \norm{h}_{\psi_2}^2 \max\{d, n^{1/2}\} \]
For every subgaussian random variable that doesn't depend on $d$ or $n$, we get a reasonable bound. Otherwise, there is another iterated inequality
we could apply called the Iterated strong Rosenthal inequality that removes this prefactor, which we ignore for now.
Applying this same strategy for general $q$ matches 
the state-of-the-art bound of $\mathbb{E} \norm{Y} \approx_q \max\{d^{q/2}, n^{1/2}\}$.
To give a sketch of this proof, assume that we can remove the prefactor in NCK. Then, we will
sketch why $\sigma(\mathcal{A}) \le \max\{d^{q/2}, n^{1/2}\}$ and $v(\mathcal{A}) \le d^{\frac{q - 1}{2}}$.

\textbf{Proof Sketch.}
Note in this setting, 
the $m \le q$th index is $(j_m, k)$,
the $(q + 1)$th index is $(j_1, \dots, j_q)$ and the $(q + 2)$th index is $k$.
One can compute
\begin{enumerate}
    \item For $\sigma$-flattening $R = \{q + 1\}, C = [q + 2] \setminus R$ we have  $\norm{\mathcal{A}_{[R \mid C]}}^2 = d^q$.
    \item For $\sigma$-flattening $C = \{q + 2\}, R = [q + 2] \setminus C$,  $\norm{\mathcal{A}_{[R \mid C]}}^2 = n$.
\end{enumerate}
For any other $\sigma$-flattening $\mathcal{A}_{[R \mid C]}$, there are $t, t' \in [q]$ such that $t \in R$
and $t' \in C$. Hence, both the indices $j_{t'}$ and $k$ appear at least once both transposed and non-transposed:
for instance $e_{j_{t'}}$ comes from the $q + 1$ index in $R$ while $e_{j_{t'}}^T$ comes from the $t'$ index in $C$.
Hence,
the norm of the $t'$ and $k$ terms squared is exactly $1$ (this is the identity case), and the norm squared of any other type of terms is at most $d$.
There are only at most $q - 1$ types besides these, hence can only have squared norm at most $d^{q - 1}$.

Similarly, for any $v$-flattening, since $R$ is non-empty, there must be some $t \in R$
with $t \le q$, which means $e_{j_t}\otimes e_k$ appears; since $q + 1, q+2$ appears in $C$,
this means $e_{j_t}^T \otimes e_k^T$ appears, so again only $q - 1$ types of terms can contribute at most $d$ to the norm.


\subsection{Nearly Combinatorial Type}
It turns out, many natural settings also are ``close'' to combinatorial type, in the sense
that they have weights on the different matrices (or distribution of the random variables multiplying the terms get scaled).
An example comes from the analysis of sum-of-squares algorithms for tensor PCA \cite{hopkinssos}. In particular, the key step
is getting a bound of the following form, which gives an upper bound on sum-of-squares algorithm performance.
\begin{theorem} Let $W_1, \dots, W_n$ be iid $d \times d$ random matrices with iid $N(0, 1)$ entries. Then,
\[ \mathbb{E} \norm{\sum_{i \in [n]} (W_i \otimes W_i - \mathbb{E}[W_i \otimes W_i])} \lesssim d\sqrt{n} \]
provided that $n, d \gtrsim \log(d + n)^4$
\end{theorem}
Note that in the original paper, there is actually an extra $\sqrt{\log n}$ factor, which our calculation will elide.

\textbf{Proof Sketch.} Define $g_{ijk} = W_i[j,k]$.
Note that we can decompose $X$ as
\[ \sum_{i \in [n], j, k \in [d]} (g_{ijk}^2 - 1) e_{(j, j)} \otimes e_{(k,k)}^T + \sum_{\substack{i \in [n]\\j_1,k_1,j_2,k_2 \in [d]}} 1_{(j_1, k_1) \neq (j_2, k_2)} g_{i, j_1, k_1} g_{i, j_2, k_2} e_{(j_1, j_2)} \otimes e_{(k_1, k_2)}^T \]
The first term can be handled with the known matrix concentration inequalities after noticing that $e_{(j,j)} \simeq e_j$
(though when taking $h_{i,j,k} = (g_{ijk}^2 - 1)$,
iterated Rosenthal will work to show an upper bound of $\tilde{O}(\sqrt{dn})$), so we focus our attention to
the second term. As alluded to before, we can pass to the decoupled version
\[Y = \sum_{\substack{i \in [n]\\j_1, k_1, j_2, k_2 \in [d] }} 1_{(j_1, k_1) \neq (j_2, k_2)} g_{i, j_1, k_1}^{(1)} g_{i,j_2,k_2}^{(2)} e_{(j_1, j_2)} \otimes e_{(k_1, k_2)}^T\]
while only gaining an error which is a constant (see that $q = 2$).

Now, imagine a version of $Y$ where there was no indicator function in front. Then one can show (again by direct computation)
that $\sigma(\mathcal{A}) = d\sqrt{n}$ and $v(\mathcal{A}) = \max\{d, \sqrt{n}\}$. Using iterated strong NCK provides
$\mathbb{E} \norm{Y} \lesssim d\sqrt{n} + \log(d^2 + nd^2)^2 \max\{d, \sqrt{n}\}$,
which would imply the proposition.

To show that nothing is lost by taking this non-indicator version, first we start with a quick lemma.
\begin{lemma}
    If $M, M'$ are real matrices so that $|M_{i,j}| \le M_{i,j}'$ for all $i, j$,
    then $\norm{M} \le \norm{M'}$.
\end{lemma}

\begin{proof*}
    $\norm{M} = \sup_{\norm{x} = \norm{y} = 1} \sum_{i,j} x_i M_{i,j} y_j \le \sup_{\norm{x} = \norm{y} = 1} \sum_{i,j} |x_i| |M_{i,j}| |y_j| \le \sup_{\norm{x} = \norm{y} = 1} \sum_{i,j} |x_i| M_{i,j}' |y_j| \le  \norm{M'}$.
\end{proof*}

Let's take a look at the flattenings of this new matrix model. By again combining terms appropriately
as we did in the Khatri-Rao case, we get that:
\[ \mathcal{A}_{[R \mid C]} \simeq \sum_{\substack{s_1=i \in [n]\\s_2=j_1, s_3=k_1 \in [d] \\ s_4=j_2, s_5=k_2 \in [d] }} 1_{(j_1, k_1) \neq (j_2, k_2)}  \bigotimes_{u \in [5]} (e_{s_u}^{\otimes \mu_u} \otimes (e_{s_u}^T)^{\otimes \nu_u}) \]
Now consider the matrix
\[ \mathcal{A}_{[R \mid C]}^{\text{ideal}} \simeq \sum_{\substack{s_1=i \in [n]\\s_2=j_1, s_3=k_1 \in [d] \\ s_4=j_2, s_5=k_2 \in [d] }} \bigotimes_{u \in [5]} (e_{s_u}^{\otimes \mu_u} \otimes (e_{s_u}^T)^{\otimes \nu_u}) \]
Then, now just note that the entries of the original matrix (which are all nonnegative) are at most the latter matrix,
so $\norm{\mathcal{A}_{[R \mid C]}} \le \norm{\mathcal{A}_{[R \mid C]}^{\text{ideal}}}$, meaning our idealization was correct.

In fact, this trick generalizes for arbitrary ``weight functions''; the penalty on the
norm of the flattening is in the infinity norm of the function.

\bibliographystyle{alpha}
\bibliography{refs}
\end{document}